\chapter{CONCLUSION}

\section{Summary of Contributions}

This thesis designed, implemented and empirically validated an end-to-end
automated competitive-intelligence pipeline for the Amazon
consumer-electronics marketplace. The contributions span four layers ---
data engineering, analytics, evaluation, and delivery:

\begin{enumerate}
  \item \textbf{Data engineering.} A production pipeline collecting daily
  Top-50 category rankings, product-level snapshots for 30 watchlist
  ASINs, and weekly customer reviews. The pipeline runs unattended on
  GitHub Actions cron (Apify $\to$ Supabase) and achieved
  $\approx$ 90\% completeness on category rankings and $\approx$ 100\%
  on watchlist snapshots over the 13-day pilot window.

  \item \textbf{Quantitative competitive metrics.} Two composite scores
  --- the Listing Quality Score (LQS) and Brand Momentum Score (BMS) ---
  give practitioners standardised, longitudinally comparable benchmarks of
  product competitiveness. K-Means price tiering provides a complementary
  segmentation view that is updated daily.

  \item \textbf{NLP-based sentiment pipeline.} A pre-trained RoBERTa
  classifier produces per-ASIN polarity scores and feeds aspect-level
  breakdowns from the Apify reviews actor. The classifier was benchmarked
  by distant supervision against star ratings on 3{,}851 reviews,
  achieving 81.2\% accuracy and macro-F\textsubscript{1} 0.594.

  \item \textbf{Probabilistic price forecasting.} A Facebook Prophet
  pipeline produces 7-day forecasts with 80\% prediction intervals,
  benchmarked under a walk-forward backtest. The forecaster matches a naive
  last-value baseline on point accuracy under the 13-day data regime
  (MAE \$0.601, MAPE 1.16\%) but provides calibrated PIs (90.2\% coverage)
  that the rule engine and downstream agents can exploit.

  \item \textbf{Operational alerting.} A rule-based engine translates
  metric changes into Critical / Action / Watch tiers; over the pilot
  window it produced 1{,}248 alerts of which 55 were high-severity.

  \item \textbf{Multi-agent delivery layer.} On top of the analytics
  warehouse, three specialist Telegram bots (Competitor Spy, Sentiment
  Detective, Momentum Strategist) backed by the self-hosted OpenClaw
  runtime call thirteen read-only Python skills against Supabase. The
  layer delivers the same intelligence as the Streamlit dashboard but
  via a push-and-pull conversational interface, including three
  scheduled cron briefs in Asia/Ho\_Chi\_Minh time.
\end{enumerate}

\section{Key Findings}

The most defensible empirical findings, as reported in detail in
Chapter~\ref{chap:results}, are:

\begin{itemize}
  \item \textbf{Three competitive regimes.} Mature/stable (TWS),
  mid-mature/brand-differentiated (Gaming Keyboards) and
  commoditised/high-churn (Portable Chargers) emerge clearly from a
  combined analysis of price-tier shape, entrant/exit churn, review
  volumes, and sentiment polarity.
  \item \textbf{Listing quality has converged.} Top-50 watchlist LQS
  scores in mature consumer-electronics categories cluster tightly between
  79.5 and 96.9 (overall mean 93.4). The metric discriminates most
  strongly at the long tail and at very recent flagship launches that have
  not yet accumulated reviews.
  \item \textbf{High-rating, mid-sentiment commodity gap.} Portable
  Chargers post a 4.48 mean star rating but only 0.294 mean RoBERTa
  sentiment, against 0.564 for TWS at a slightly lower 4.33 stars. Star
  averages can therefore mask substantial customer concern in the
  free-text content --- the precise gap that motivates NLP-based
  sentiment monitoring on top of star tracking.
  \item \textbf{Top-50 image stability.} Zero pHash-detected hero-image
  changes and zero \texttt{bsr\_drop} alerts fired during the 13-day
  window, suggesting that competitive moves at the top of mature
  electronics categories operate primarily through pricing, ad share, and
  rank churn rather than visual A/B testing on a daily cadence.
  \item \textbf{Sentiment pipeline is fit-for-purpose.} The RoBERTa
  classifier achieves strong positive (F\textsubscript{1} 0.908) and
  negative (0.715) discrimination on noisy distant labels; the
  weak neutral class (0.158) is consistent with the inherent ambiguity of
  3-star reviews and does not affect the directional polarity signal that
  the BMS and LQS metrics consume.
  \item \textbf{Forecast under data scarcity.} On a 13-day window
  Prophet does not beat a naive last-value baseline on point accuracy,
  but its 80\% prediction interval covers 90.2\% of hold-out actuals ---
  a calibrated probabilistic envelope is the deliverable, not a
  point prediction.
\end{itemize}

\section{Research Limitations}

\begin{enumerate}
  \item \textbf{Observation window.} A 13-day window is too short for
  seasonality, holiday-effect, or longitudinal-trend analysis. Findings
  should be considered illustrative rather than statistically conclusive.

  \item \textbf{Category scope.} Three consumer-electronics
  sub-categories may not generalise to broader marketplace dynamics ---
  in particular the commoditisation finding for Portable Chargers is
  category-specific and would not be expected in luxury or apparel
  verticals.

  \item \textbf{Distant supervision.} The RoBERTa accuracy of 81.2\% is
  bounded above by the noisiness of star ratings as a sentiment proxy.
  The collapse on the neutral class reflects both genuine model
  uncertainty and the inherent ambiguity of the truth label.

  \item \textbf{Forecast horizon under short history.} Prophet's
  inability to outperform the naive baseline on a 3-day hold-out is
  expected given a $\le 13$-day training window in stable consumer
  categories; longer histories or richer regressors (promotion calendars,
  competitor price moves) would shift the comparison.

  \item \textbf{Causal inference.} All analyses in this thesis are
  correlational. Claims about the relationship between LQS, sentiment,
  and rank performance are descriptive associations and cannot be
  interpreted causally without a controlled experiment.

  \item \textbf{Operational fragility.} The OpenClaw Gateway runs in a
  self-managed VMware Ubuntu VM whose NAT layer occasionally drops the
  long-poll connection to Telegram. Recovery is automatic but introduces
  30--60s message lags. A managed cloud-VM deployment would remove this
  artefact at the cost of operational complexity.
\end{enumerate}

\section{Future Work}

Several directions naturally extend this research:

\begin{itemize}
  \item \textbf{Extended collection.} Running the pipeline for 3--6
  months would enable seasonality analysis (Prime Day, Black Friday,
  Cyber Monday), would let the Prophet forecaster compete against the
  naive baseline more meaningfully, and would allow drift detection in
  the LQS / BMS distributions.

  \item \textbf{Domain-adaptive sentiment.} Fine-tuning RoBERTa on a
  small gold-labelled subset of Amazon electronics reviews would close the
  reported gap on the negative and (especially) neutral classes, and
  would make per-aspect classification more reliable.

  \item \textbf{Sponsored-share integration.} Bringing daily share-of-voice
  in sponsored slots into the BMS as a fourth term would broaden the
  metric beyond organic momentum into paid competitive pressure.

  \item \textbf{Cross-marketplace replication.} Extending the Apify
  collectors to Amazon UK, DE, and JP would provide a comparative view
  of pricing and listing-quality strategy across geographies, and would
  test whether the three-regime competitive taxonomy (mature / mid /
  commoditised) generalises across markets.

  \item \textbf{Agent self-evaluation.} The OpenClaw Gateway already
  records per-skill latency and per-message LLM token usage. A
  natural extension is for each agent to self-score the relevance and
  freshness of its own briefs (e.g.\ via a separate evaluator prompt) and
  surface low-confidence answers to the user explicitly.
\end{itemize}
